\section{El taller a la luz de la teoría}
\begin{center}\small\renewcommand{\arraystretch}{1.3}
\begin{longtable}{@{}p{2.3cm}p{4.3cm}p{8.8cm}@{}}\toprule
\textbf{Ejercicio} & \textbf{Qué evalúa} & \textbf{Clave y trampa}\\ \midrule\endhead
I y II, verificar & Axiomas de topología & Tablas de $\cup$ y $\cap$ de los pares no comparables. $\tau_2$ falla una
unión, $\tau_3$ una intersección. En II hay una errata en $\tau_1$.\\
I(a), II(a) & Cerrados & Complemento de cada abierto. Hay conjuntos abiertos y cerrados a la vez.\\
I(b), II(b) & Derivado & Usar el abierto mínimo y quitar el propio punto. Revisar \emph{todos} los puntos de $X$.\\
I(c), II(c) & Clausura & Menor cerrado que contiene al conjunto; comprobar con $A\cup A'$.\\
I(d), II(d) & Topología relativa & Intersecar cada abierto con $H$ y quitar repetidos.\\
I(e) & Densidad & $\cl{U}=X$. No confundir «denso» con «grande».\\
I(f), II(e) & Vecindades & Todos los superconjuntos de $U_x$. Convención de vecindad no necesariamente abierta.\\
I(g), II(f) & Interior & Unión de los abiertos contenidos. Puede ser $\vac$ o muy pequeño.\\
I(h), II(g) & Frontera ($\sigma$) & $\cl{A}\sm\inte{A}$. Notación por confirmar.\\
III.1--3 & Topologías pequeñas & Criterio de cadena o complementarios (4 términos); tres formas (5 términos).\\
III.4 & Comparación & Inclusión entre topologías; mismas cantidades de abiertos $\Rightarrow$ no comparables.\\
III.5 & Axiomas & Igual que la Parte I, con seis puntos.\\
III.6, III.7 & Topologías en conjuntos infinitos & Demostrar (T1)--(T3) con elementos; el abierto mínimo da las
vecindades. En III.7 falta $\vac$ en la definición.\\
III.8, III.9 & Topología generada & Primero intersecciones finitas, después uniones. En III.9 la sub-base no cubre $X$.\\
III.10 & Sub-bases & Familias que generan exactamente la topología dada.\\
III.11 & Sub-base en $\R$ & $[x-1,x]\cap[x,x+1]=\{x\}$: sale la discreta.\\
III.12 & Bases & Agrandar una base con abiertos sigue dando base. Errata en la conclusión.\\
III.13 & Axiomas & $A\subseteq B$, $B\subseteq A$ o $B=X\sm A$.\\
IV.1 & Continuidad global y puntual & Imágenes inversas; continua en $d$ pero no globalmente.\\
IV.2 & Función abierta & Contraejemplo $(-1,1)\mapsto[0,1)$.\\
IV.3 & Continua, abierta, cerrada & Tres pruebas independientes por función. $\tau_1$ no es topología; $X\not\homeo Y$.\\
IV.4 & Homeomorfismo & Contar abiertos para negar; construir la biyección para afirmar.\\ \bottomrule
\end{longtable}
\end{center}

% =====================================================================
\section{Recetario para el parcial}\label{sec:recetario}
\begin{center}\small\renewcommand{\arraystretch}{1.35}
\begin{longtable}{@{}p{4.3cm}p{11.1cm}@{}}\toprule
\textbf{Si piden\dots} & \textbf{Se hace\dots}\\ \midrule\endhead
¿Es topología? & (1) ¿Están $\vac$ y $X$? (2) Uniones e intersecciones de los pares no comparables. Un fallo basta
para decir que no.\\
Cerrados & Complemento de cada abierto.\\
Abierto mínimo $U_x$ & Intersección de los abiertos que contienen a $x$. Hacer la tabla al principio: sirve para todo
lo demás.\\
Vecindades $\N_x$ & Todos los conjuntos que contienen a $U_x$.\\
Derivado $A'$ & $x\in A'\iff(U_x\sm\{x\})\cap A\neq\vac$, para cada $x\in X$.\\
Clausura $\cl{A}$ & Menor cerrado que contiene a $A$; o $A\cup A'$; o $\{x:U_x\cap A\neq\vac\}$.\\
Interior $\inte{A}$ & Unión de los abiertos contenidos en $A$; o $\{x:U_x\subseteq A\}$.\\
Exterior & $\Int(X\sm A)=X\sm\cl{A}$.\\
Frontera & $\cl{A}\sm\inte{A}$. Comprobar $X=\inte{A}\sqcup\Fr(A)\sqcup\Ext(A)$.\\
¿Denso? & $\cl{D}=X$; o $D$ corta a todo abierto no vacío.\\
Topología relativa $\tau_H$ & $\{U\cap H:U\in\tau\}$, sin repetidos.\\
¿Comparables? & ¿Una lista de abiertos está contenida en la otra?\\
Topología generada & Intersecciones finitas (más $X$) $\to$ base; uniones (más $\vac$) $\to$ topología; verificar.\\
¿Base de $\tau$? & Miembros abiertos y todo abierto es unión de ellos.\\
¿Continua? & $f^{-1}(V)$ abierto en la salida para cada abierto $V$ de la llegada.\\
¿Continua en $p$? & Para cada abierto $V\ni f(p)$, ¿$f^{-1}(V)$ es vecindad de $p$? Finito: $f(U_p)\subseteq V_{f(p)}$.\\
¿Abierta? & $f(U)$ abierto en la llegada para cada abierto $U$ de la salida. Descarte: $f(X)$.\\
¿Cerrada? & $f(F)$ cerrado en la llegada para cada cerrado $F$ de la salida. Descarte: $f(X)$.\\
¿Homeomorfos? & Negar con un invariante (número de abiertos, de puntos abiertos\dots) o construir la biyección y
tabular $U\mapsto\varphi(U)$.\\
Demostrar que algo es topología (caso infinito) & (T1) mostrar que $\vac$ y $X$ cumplen la condición; (T2) tomar una
familia arbitraria $\{A_i\}$ y un elemento de la unión; (T3) tomar $A$, $B$ y un elemento de la intersección.\\ \bottomrule
\end{longtable}
\end{center}

% =====================================================================
\section{Errores frecuentes}
\begin{enumerate}[leftmargin=1.8em,itemsep=0.35em]
\item \textbf{Verificar solo uniones o solo intersecciones.} Hacen falta las dos (ejemplo: $\tau_3$ de la Parte~I).
\item \textbf{Pensar que cerrado es «no abierto».} Hay conjuntos con las dos propiedades y conjuntos sin ninguna.
\item \textbf{Olvidar el «distinto de $x$» en el derivado.} Un punto de $A$ puede no estar en $A'$.
\item \textbf{Buscar puntos de acumulación solo dentro de $A$.} Hay que revisar todo $X$.
\item \textbf{Confundir $A'$ con $\cl{A}$.} $\cl{A}=A\cup A'$; el derivado puede ser más pequeño o no contener a $A$.
\item \textbf{Creer que la frontera es «pequeña».} En espacios finitos puede ser casi todo $X$.
\item \textbf{Dar como $\N_x$ solo los abiertos} cuando la definición del curso admite vecindades no abiertas (o al
revés). Aclarar la convención en la respuesta.
\item \textbf{En la topología generada, unir antes de intersecar}, u olvidar agregar $X$ y $\vac$.
\item \textbf{Para la continuidad, calcular $f(U)$} en vez de $f^{-1}(V)$.
\item \textbf{Concluir «homeomorfos» porque hay una biyección} o porque coincide el número de abiertos. Hay que
exhibir la biyección y comprobar que lleva abiertos en abiertos en ambos sentidos.
\item \textbf{Decir «abierto» sin decir en qué espacio}, cuando hay subespacios de por medio.
\item \textbf{No leer el enunciado con cuidado.} Este taller tiene al menos cinco erratas; en un examen, si un dato
parece inconsistente, hay que señalarlo y explicar qué se asumió.
\end{enumerate}

% =====================================================================
\section{Autoevaluación (con respuestas)}\label{sec:auto}
Sea $X=\s{p,q,r,s}$ y $\tau=\{\vac,\s{p},\s{p,q},\s{p,r},\s{p,q,r},X\}$.
\begin{enumerate}[leftmargin=1.8em,itemsep=0.25em]
\item Compruebe que $\tau$ es una topología.
\item Halle los cerrados.
\item Halle el abierto mínimo de cada punto.
\item Para $A=\s{q,s}$ halle $A'$ y $\cl{A}$.
\item Halle $\cl{\s{p}}$ y $\s{p}'$. ¿Es $\s{p}$ denso?
\item Halle el interior de $\s{q,r,s}$ y el de $\s{p,q,s}$.
\item Para $B=\s{p,s}$ halle interior, clausura, frontera y exterior.
\item Halle $\N_s$ y $\N_q$.
\item Halle la topología relativa sobre $H=\s{q,r,s}$.
\item ¿Cuáles subconjuntos de $X$ son densos?
\item Sea $f:X\to X$ con $f(p)=p$, $f(q)=p$, $f(r)=q$, $f(s)=r$. ¿Es continua? ¿Abierta? ¿Cerrada?
\item Sea $\varphi:X\to X$ la función que intercambia $q$ con $r$ y fija $p$ y $s$. ¿Es un homeomorfismo?
\item Halle la topología $\tau^{\ast}$ sobre $X$ generada por $\C=\{\s{p,q},\s{q,r}\}$.
\item ¿Son homeomorfos $(X,\tau)$ y $(X,\tau^{\ast})$?
\item Verdadero o falso: (a) en un espacio discreto todo subconjunto es abierto y cerrado; (b) si $A$ es abierto,
entonces $A'=\vac$; (c) $\cl{A\cap B}=\cl{A}\cap\cl{B}$ siempre; (d) todo homeomorfismo es una función abierta; (e) la
unión de dos topologías es una topología; (f) toda biyección continua es homeomorfismo; (g) $\Fr(A)=\Fr(X\sm A)$
siempre; (h) el interior de un conjunto no vacío es no vacío.
\end{enumerate}

\begin{respuesta}
\textbf{\sffamily\color{verde!70!black}Respuestas.}
\begin{enumerate}[leftmargin=1.8em,itemsep=0.2em]
\item El único par no comparable es $\s{p,q}$, $\s{p,r}$: unión $\s{p,q,r}$, intersección $\s{p}$; ambos están.
\item $\cotau=\{\vac,\s{s},\s{q,s},\s{r,s},\s{q,r,s},X\}$.
\item $U_p=\s{p}$, $U_q=\s{p,q}$, $U_r=\s{p,r}$, $U_s=X$.
\item $A'=\s{s}$ (solo $U_s\sm\s{s}=\s{p,q,r}$ corta a $A$); $\cl{A}=\s{q,s}$, que es cerrado.
\item $\cl{\s{p}}=X$ y $\s{p}'=\s{q,r,s}$. Sí es denso.
\item $\Int\s{q,r,s}=\vac$; $\Int\s{p,q,s}=\s{p,q}$.
\item $\inte{B}=\s{p}$, $\cl{B}=X$, $\Fr(B)=\s{q,r,s}$, $\Ext(B)=\vac$.
\item $\N_s=\{X\}$; $\N_q=\{\s{p,q},\s{p,q,r},\s{p,q,s},X\}$.
\item $\tau_H=\{\vac,\s{q},\s{r},\s{q,r},H\}$.
\item Exactamente los que contienen a $p$ (8 conjuntos): todo abierto no vacío contiene a $p$.
\item No es continua: $f^{-1}(\s{p,r})=\s{p,q,s}\notin\tau$. Sí es abierta: las imágenes de los abiertos son
$\s{p},\s{p},\s{p,q},\s{p,q},\s{p,q,r}$. No es cerrada: $f(\s{s})=\s{r}$ no es cerrado.
\item Sí: es biyectiva y lleva $\tau$ sobre $\tau$ (intercambia $\s{p,q}$ con $\s{p,r}$ y fija los demás abiertos).
\item Base: $\{X,\s{q},\s{p,q},\s{q,r}\}$. $\tau^{\ast}=\{\vac,\s{q},\s{p,q},\s{q,r},\s{p,q,r},X\}$.
\item Sí: $\psi=(p\mapsto q,\ q\mapsto p,\ r\mapsto r,\ s\mapsto s)$ lleva $\s{p}\mapsto\s{q}$, $\s{p,q}\mapsto\s{p,q}$,
$\s{p,r}\mapsto\s{q,r}$, $\s{p,q,r}\mapsto\s{p,q,r}$, que es exactamente $\tau^{\ast}$.
\item (a) V. (b) F: en el ejercicio 5, $\s{p}$ es abierto y $\s{p}'\neq\vac$. (c) F: solo vale «$\subseteq$». (d) V.
(e) F. (f) F: la inversa puede no ser continua. (g) V. (h) F: ejercicio 6.
\end{enumerate}
\end{respuesta}

% =====================================================================
\section{Glosario de notación}
\begin{center}\small\renewcommand{\arraystretch}{1.3}
\begin{tabular}{@{}p{3.0cm}p{5.2cm}p{7.2cm}@{}}\toprule
\textbf{Aquí} & \textbf{Significado} & \textbf{Otras notaciones que aparecen en los textos}\\ \midrule
$\tau$ & topología (familia de abiertos) & $\mathcal{T}$, $\mathcal{O}$\\
$\cotau$ & familia de cerrados & $c_0\tau$ (taller), $\mathcal{F}$, $\mathcal{C}$\\
$\N_x$ & sistema de vecindades de $x$ & $N_x$ (taller), $\mathcal{V}(x)$, $\mathcal{U}_x$, $\mathcal{E}(x)$ (entornos)\\
$U_x$ & abierto mínimo de $x$ (espacios finitos) & $\mathrm{st}(x)$, $\uparrow\! x$\\
$A'$ & derivado & $A^{d}$, $\mathrm{Der}(A)$, $\mathrm{ac}(A)$\\
$\cl{A}$ & clausura o adherencia & $\mathrm{cl}(A)$, $\mathrm{Cl}\,A$, $A^{-}$, $\mathrm{adh}(A)$\\
$\inte{A}$, $\Int(A)$ & interior & $A^{\circ}$, $\mathrm{Int}\,A$\\
$\Ext(A)$ & exterior & $\mathrm{Ext}(A)$\\
$\Fr(A)$ & frontera & $\partial A$, $b(A)$, $\mathrm{Bd}\,A$; $\sigma(A)$ en el taller (por confirmar)\\
$\tau_H$ & topología relativa a $H$ & $\tau|_H$, $\tau/H$\\
$\B$, $\Sb$ & base, sub-base & $\mathfrak{B}$, $\mathcal{S}$\\
$\tau(\C)$ & topología generada por $\C$ & $\langle\C\rangle$\\
$X\homeo Y$ & $X$ homeomorfo a $Y$ & $X\approx Y$, $X\simeq Y$\\ \bottomrule
\end{tabular}\end{center}

% =====================================================================
\section*{Fuentes}
\addcontentsline{toc}{section}{Fuentes}
\markboth{Fuentes}{}
\begin{ojo}[Alcance de la verificación de fuentes]
El entorno en que se preparó esta guía permitía \emph{buscar} en la web, pero no abrir directamente las páginas ni los
PDF. Por eso: (1) los recursos en línea se consultaron a través de los resultados y resúmenes de búsqueda, no leyendo
el documento completo; (2) los libros impresos se citan por capítulo o sección según su índice, sin número de página,
y conviene confirmar la numeración en la edición que se tenga a mano; (3) todo resultado matemático de la guía se
puede comprobar de forma independiente: las demostraciones están escritas y los cálculos finitos se verificaron por
computadora. No se encontraron apuntes públicos del curso de la Universidad de Panamá, de modo que las convenciones de
notación propias del curso ($c_0\tau$, $\sigma$) quedan como supuestos.
\end{ojo}
\textbf{Libros de texto}
\begin{enumerate}[label={[\arabic*]},leftmargin=2.6em,itemsep=0.3em]\small\sloppy
\item S.~Lipschutz, \emph{Theory and Problems of General Topology}, Schaum's Outline Series, McGraw-Hill, Nueva York,
1965 (ed. en español: \emph{Topología general}). Capítulos «Topological spaces: definitions», «Bases and subbases» y
«Continuity and topological equivalence» (5, 6 y 7). Es, con mucha probabilidad, la fuente de los ejemplos de las
Partes I y IV del taller; la topología $\{X,\vac,\{a\},\{c,d\},\{a,c,d\},\{b,c,d,e\}\}$ es su ejemplo recurrente.
\item J.~R.~Munkres, \emph{Topology}, 2.ª ed., Prentice Hall, 2000. Capítulo 2, «Topological Spaces and Continuous
Functions»: §12 (espacios topológicos), §13 (bases), §16 (subespacios), §17 (cerrados y puntos límite), §18
(funciones continuas).
\item J.~L.~Kelley, \emph{General Topology}, Van Nostrand, 1955; reimpr. Springer, Graduate Texts in Mathematics 27.
Capítulo 1.
\item S.~Willard, \emph{General Topology}, Addison-Wesley, 1970; reimpr. Dover, 2004. Capítulos 2 y 3.
\item R.~Engelking, \emph{General Topology}, ed. revisada, Heldermann, Berlín, 1989. Capítulo 1.
\item G.~N.~Rubiano, \emph{Topología general (un primer curso)}, 2.ª ed., Universidad Nacional de Colombia, Bogotá,
2002.
\item M.~Macho Stadler, \emph{Topología}, apuntes del curso 2014--2015, Universidad del País Vasco (UPV/EHU).
\url{https://www.ehu.eus/~mtwmastm/Topologia1415.pdf}
\item F.~Cano, M.~Martín-Vega y B.~Molina-Samper, \emph{Topología: notas de un curso}, Universidad de Valladolid.
\url{https://uvadoc.uva.es/handle/10324/70838}
\item S.~A.~Morris, \emph{Topology Without Tears}, libro de acceso libre. \url{https://www.topologywithouttears.net}
\item L.~A.~Steen y J.~A.~Seebach, \emph{Counterexamples in Topology}, 2.ª ed., Springer, 1978. Ejemplos 5 (topología
de partición), 6 (par--impar) y 8 a 10 (punto particular), según los listados consultados.
\end{enumerate}
\textbf{Recursos en línea y artículos}
\begin{enumerate}[label={[\arabic*]},leftmargin=2.6em,itemsep=0.3em,start=11]\small\sloppy
\item OEIS Foundation, A000798: \emph{Number of different quasi-orders (or topologies, or transitive digraphs) with
$n$ labeled elements}. \url{https://oeis.org/A000798}
\item OEIS Foundation, A001930: \emph{Number of topologies, or transitive digraphs, with $n$ unlabeled nodes}.
\url{https://oeis.org/A001930}
\item Wikipedia, «Finite topological space». \url{https://en.wikipedia.org/wiki/Finite_topological_space}
\item Wikipedia, «Alexandrov topology». \url{https://en.wikipedia.org/wiki/Alexandrov_topology}
\item Wikipedia, «Kuratowski closure axioms». \url{https://en.wikipedia.org/wiki/Kuratowski_closure_axioms}
\item Wikipedia, «Subbase». \url{https://en.wikipedia.org/wiki/Subbase}
\item Wikipedia, «Comparison of topologies». \url{https://en.wikipedia.org/wiki/Comparison_of_topologies}
\item Wikipedia, «Particular point topology». \url{https://en.wikipedia.org/wiki/Particular_point_topology}
\item ProofWiki, «Definition: Odd-Even Topology». \url{https://proofwiki.org/wiki/Definition:Odd-Even_Topology}
\item B.~G.~Wachsmuth, \emph{Interactive Real Analysis}, §6.4, ejemplo de $x^2$ y la imagen de $(-1,1)$.
\url{https://mathcs.org/analysis/reals/cont/answers/imag1.html}
\item M.~Benoumhani, «The number of topologies on a finite set», \emph{Journal of Integer Sequences} 9 (2006),
artículo 06.2.6.
\item J.~Evans, \emph{Topology and Groups} (University College London), «2.07 Homeomorphisms».
\url{https://www.homepages.ucl.ac.uk/~ucahjde/tg/html/topsp07.html}
\item R.~Tinarrage, \emph{Topological Data Analysis} (EMAp/FGV), lección 2: homeomorfismos.
\url{https://raphaeltinarrage.github.io/files/EMAp/Lesson2.pdf}
\item Universidad de Extremadura, apuntes de \emph{Fundamentos Matemáticos}, tema 3 (notación $\mathrm{Fr}(A)$ y
$\mathrm{Ext}(A)$). \url{https://matematicas.unex.es/~trinidad/fmi/TTema3.pdf}
\item R.~P.~Stanley, «On the number of open sets of finite topologies», \emph{Journal of Combinatorial Theory} 10
(1971), 74--79.
\item N.~Bourbaki, \emph{Topologie générale}, capítulo 1 (estructuras topológicas).
\item K.~Ragnarsson y B.~E.~Tenner, «Obtainable sizes of topologies on finite sets», \emph{Journal of Combinatorial
Theory, Series A} 117 (2010), 138--151; arXiv:0802.2550.
\end{enumerate}
